\section{Introduction}\label{sec:intro}

\IEEEPARstart{C}{ylinder} chess is ordinary chess on a board whose a-file and h-file are joined
\cite{Pritchard}: a white pawn on h4 can capture on a5, a king on a1 can step to h2, and there are no side edges or
corners. We use combinatorial game theory (CGT) to study its king-and-pawn endgames.

Under normal play the player who cannot move loses; CGT gives each position a value, and a sum of independent
games has the sum of their values \cite{ONAG,WinningWays,Siegel2013}. Chess can be drawn and its kings roam,
yet Elkies \cite{Elkies1996} showed that many pawn endgames split into local fights with exact values, under
two conventions:
\begin{itemize}
\item a passed pawn wins, and is worth a large number $\pm\BIG$;
\item the kings stand in a mutual zugzwang, the \emph{tr\'ebuchet}, worth $0$.
\end{itemize}

The cylinder changes this picture in three ways. There are no rook pawns, and shifting a position sideways
does not change its value. A position with a pawn on every file, a \emph{ring}, has no edge at which to cut
it. And the corners that decide mates and fortresses are gone.

A nimber $\nim{k}$ is a game in which either player may move to any of $\nim{0}=0$, $\nim{1}=\star$,
\dots, $\nim{(k-1)}$. The simplest after $\star=\game{0}{0}$ is $\nim{2}=\game{0,\star}{0,\star}$: either player
may move to $0$ or to $\star$. Elkies built higher nimbers on large boards from Dawson's pawn game \cite{Elkies2002,Dawson1935}, and
Drummond-Cole built positions of value $\nim{2}$ in chess with the help of pieces other than kings and
pawns~\cite{DrummondCole2005}. We ask whether $\nim{2}$ occurs in a king-and-pawn endgame on an 8-rank board,
the \emph{$\nim{2}$ question}, and whether real chess agrees.

This paper extends an unpublished 2025 study of the basic values\bytheauthor~\cite{Study2025}: we check
every position by computer, correct its errors and carry the analysis on, in the position-by-position format
of Wu, Y\"u and Landry~\cite{WuYuLandry2010}. Theorems and headline numbers
carry labels: \lab{P}, proved here; \lab{AI}, proof given or sketched here (in full in the
repository), checked by AI models, human check pending; \lab{C2}, two independently written programs agree,
the second written from the rules alone; \lab{C1}, one program; \lab{Conj}, a conjecture. Our contributions
are:
\begin{itemize}
\item \emph{Model} (Sections~\ref{sec:rules} and~\ref{sec:sums}): a passed-pawn test that respects en passant,
  a cut lemma \lab{P}, a sum-safety theorem \lab{AI}, and a position showing that checking only the
  top-level components is not enough \lab{C2}.
\item \emph{Kings} (Section~\ref{sec:kings}): conditions (V1)--(V3) under which a frozen tr\'ebuchet gives the
  result of the whole-board pawn game \lab{AI}, and (V4) for the $\pm\BIG$ value \lab{AI}.
\item \emph{Small values} (Sections~\ref{sec:basic}, \ref{sec:atlas} and~\ref{sec:targets}): cylinder-only values
  come from rings \lab{P}; $1/8$ and $\upsec\star$ \lab{C2}, and Tiny~1 \lab{C1}, occur in 10-man positions
  with a legal pawn history.
\item \emph{The $\nim{2}$ question} (Section~\ref{sec:star2}): in the default setting (count rule, e.p.-aware
  test; Section~\ref{sec:ep}) no $\nim{2}$ has six
  pawns or fewer on ordinary boards of any width, but $\nim{2}$ appears with seven
  \lab{C2}; on the cylinder $\cyl{3}$ a six-pawn ring suffices \lab{C2}, as a whole-board game some ring with $n+3$
  pawns is $\nim{2}$ on every $\cyl{n}$ with $n\ge5$ \lab{AI}, and on $\cyl{7}$ some ten-pawn rings are
  $\nim{3}$ \lab{C2}.
\item \emph{Real chess} (Section~\ref{sec:real}): with the new test the textbook $\nim{2}$ candidate is not
  $\nim{2}$, however draws are treated; with queen promotions only and frozen kings, the oracle value of the $1/8$ position
  is exactly $1/8$ \lab{C1}.
\item \emph{Tablebases} (Section~\ref{sec:tb}): king and rook win only 2.80\% of cylinder positions, and on the
  torus only two major pieces force mate \lab{C2}.
\item \emph{Dawson's chess} (Section~\ref{sec:dawson}): the $3\times n$ cylinder is a first-player win if and
  only if the $3\times(n-3)$ strip has Grundy value $0$ \lab{P}.
\item \emph{Engines} (Section~\ref{sec:engines}): Stockfish is weakest on infinitesimals at low depth, and on sums
  of components engines often miss the model's winning moves \lab{C1}.
\end{itemize}

Sections~\ref{sec:rules} and~\ref{sec:basic} set up the rules and the basic values,
Sections~\ref{sec:kings} and~\ref{sec:sums} add kings and sums, Sections~\ref{sec:atlas}--\ref{sec:star2}
search for small values and $\nim{2}$, and Section~\ref{sec:real} tests them against real chess.
Sections~\ref{sec:tb}--\ref{sec:engines} turn to tablebases, Dawson's chess and engines,
Section~\ref{sec:verify} gives the cross-checks and the code, Section~\ref{sec:conclusion} lists open
problems, and the Appendix gives the FEN of each position with kings.
